\section{Open questions for follow-up}

\begin{enumerate}
  \item \textbf{Does pulse-efficient CR survive when the base CX is
    ML-optimised, not vendor-default?}\\
    \emph{Basis.} Linear area-scaling of the calibrated CX only
    holds if the base pulse is well approximated as
    amplitude-linear in the target rotation. Aggressively shaped
    (GRAPE / RL) pulses are non-linear in this sense.
    \emph{Next steps.} On Qiskit-Dynamics or IBM Q, substitute a
    GRAPE-optimised CX at the same duration and remeasure
    RZX($\theta$) process fidelity for $\theta\in[0.05,\pi/2]$;
    compare against linearly-scaled default CX.

  \item \textbf{Pair-to-pair and chip-to-chip transferability of a
    single PE-CR recipe.}\\
    \emph{Basis.} Paper demonstrates on one pair of one device
    (\texttt{ibmq\_mumbai}). CR Hamiltonian coefficients vary
    across pairs; C4 (``no additional calibration'') was only
    supported for the same pair.
    \emph{Next steps.} Hamiltonian tomography on 8 pairs of one
    Falcon device, extract ZX coefficient per pair, test single
    scale-factor recipe. Repeat on a Heron (tunable-coupler)
    device.

  \item \textbf{Interaction with modern noise-aware compilation
    (DD, PEC/ZNE routing, Sabre-with-noise).}\\
    \emph{Basis.} Qiskit 0.22 transpiler predates the current
    noise-aware stack; naive composition may cancel gains
    (DD re-lengthening shortened CR windows) or compound them.
    \emph{Next steps.} \{PE on/off\} $\times$ \{DD on/off\}
    $\times$ \{ZNE routing on/off\} matrix on a QAOA benchmark
    under realistic Qiskit-Dynamics noise; publish interaction
    table.

  \item \textbf{Transferability to higher-anharmonicity qubits
    (fluxonium / heavy-fluxonium).}\\
    \emph{Basis.} PE-CR logic depends on the CR effective
    Hamiltonian being ZX-dominated. Fluxonium two-qubit gates
    (microwave-activated CZ, differential ac-Stark) do not
    obviously admit the same scaling.
    \emph{Next steps.} Simulate fluxonium-fluxonium pair in
    Qiskit-Dynamics/scQubits, test whether area-scaling of the
    microwave-activated CZ pulse produces a valid continuous-angle
    RZZ/RXX; derive alternative scaling if not.

  \item \textbf{Closed-loop pulse-shape learning for the small-$\theta$
    regime where the coherence-limited model over-predicts.}\\
    \emph{Basis.} This replication's C1 sweep gives up to 90\%
    error reduction at small $\theta$; hardware caps at
    $\sim 50\%$. The gap likely reflects control crosstalk,
    dephasing, and readout error that a closed-loop shape
    optimiser could partially compensate.
    \emph{Next steps.} CMA-ES over 4--6 pulse parameters
    (Gaussian $+$ DRAG $+$ higher-order compensation), maximise
    RZX($\theta$) process fidelity for
    $\theta\in\{\pi/16,\pi/8,\pi/4\}$; report residual gap
    between coherence-limited theory and best-optimised hardware.
\end{enumerate}
